%% ---------------------------------------------------------------------
\subsection{Method 5A: one network for every size (new)}
\label{sec:5a}

\emph{Idea.} Replace the per-size sticker network of \Mfour with one that reads the puzzle as a set of cluster tokens, so the same
weights apply at every $N$, and use its per-cluster costs to guide a staged search on sizes too large for full-puzzle beam search.

\emph{Steps.}
\begin{enumerate}
\item \emph{Tokens.} Each cluster becomes one token: the one-hot colours of its stickers in the canonical order, projected by a
  linear map shared by all clusters of the same type, plus a learned embedding of the cluster kind and a small MLP of geometric
  features (the cluster's two in-face depths divided by $N-1$, and $N/20$).
\item \emph{Encoder.} A transformer encoder~\cite{vaswani2017attention} without positional order mixes the tokens, since they form a
  set~\cite{lee2019settransformer}.
\item \emph{Additive head.} Each token predicts a non-negative cost; the costs are summed and a small global correction is added. A sum
  matches the structure of the problem, since total distance grows with the number of unsolved clusters, and it extends to any
  number of tokens.
\item \emph{Training.} Fresh non-backtracking random walks on every training size ($N=4,\dots,8$), target the walk length $k$, loss on
  $(\hat k-k)/(N^2/\ln N)$ so that all sizes contribute on one scale. The model in the main table is the GPU-trained one (width 192,
  four layers, six heads, 20,000 steps of 512 examples on a Colab T4).
\item \emph{Search, $N\le7$.} Full-puzzle beam search (Algorithm~\ref{alg:m4} with this network).
\item \emph{Search, $N\ge8$: staged beam search} (Figure~\ref{fig:pipeline}, Algorithm~\ref{alg:staged}). Following Demaine et
  al.~\cite{demaine2011algorithms}, the Square is solved in stages from the outside in: the \emph{frame} (corner, center and
  edge-cross clusters), then edge clusters two at a time, then cross and interior clusters three at a time. Each stage adds a few
  pieces to an exact goal: those pieces, and every piece solved earlier, must show their solved stickers, so later stages cannot
  undo earlier work. A stage searches only the moves that touch one of its targets plus the macro actions that do: for every row $r$
  and column $c$, the commutator $c\,r\,c\,r$, which changes only the cluster at their crossing (Demaine's Lemma~2), and the
  conjugates $r\,c\,r$ and $c\,r\,c$. The stage heuristic is the sum of \FiveA's per-token costs over the stage's target clusters and
  the earlier clusters its actions can move. If the projection of a stage's goal stickers has at most $10^5$ states, it is enumerated
  and solved optimally by lookup (the frame always qualifies: for any $N$ it is a \Sq{3}). Each stage starts at width 128 and
  retries at $4\times$ the width up to 2,048.
\end{enumerate}

\begin{figure}[t]
\centering
\begin{tikzpicture}[font=\footnotesize,
  box/.style={draw=black!55, rounded corners=2pt, align=center, inner sep=2.5pt, minimum height=7mm},
  hl/.style={box, draw=blue!65!black, fill=blue!6, line width=0.7pt},
  ar/.style={-{Stealth[length=1.6mm]}, draw=black!60}, lab/.style={font=\scriptsize, text=black!70}]
\node[box, text width=21mm] (plan) at (0,0) {stage plan\\{\scriptsize frame, edges,\\crosses, interiors}};
\node[box, text width=19mm] (small) at (28mm,0) {projected stage\\$\le 10^5$ states?};
\node[box, text width=15mm] (bfs) at (54mm,0) {exact BFS:\\optimal moves};
\node[box, text width=21mm] (inp) at (0,-17mm) {\FiveA token costs\\{\scriptsize + verified\\commutators}};
\node[hl, text width=19mm] (beam) at (28mm,-17mm) {beam search\\{\scriptsize widen until found}};
\node[box, text width=19mm] (apply) at (28mm,-32mm) {apply moves};
\draw[ar] (plan) -- (small);
\draw[ar] (small) -- node[lab, above]{yes} (bfs);
\draw[ar] (small) -- node[lab, right]{no} (beam);
\draw[ar] (inp) -- (beam);
\draw[ar] (beam) -- (apply);
\draw[ar] (bfs.south) |- (apply.east);
\draw[ar] (apply.west) -- node[lab, below]{next stage} ++(-31mm,0) |- (plan.west);
\end{tikzpicture}
\caption{\FiveA's staged search for $N\ge8$. Each stage adds a few pieces to an exact goal. Tiny stages are solved optimally by
lookup; the rest by beam search over single moves and verified commutators, guided by \FiveA's per-cluster costs.}
\label{fig:pipeline}
\end{figure}

\begin{algorithm}[t]
\caption{Staged beam search (\FiveA for $N\ge8$)}
\label{alg:staged}
\small
\KwIn{state $s$, stage plan $S_1,\dots,S_m$, cost model $h$}
\KwOut{move sequence $\mu$ with $s\cdot\mu$ solved}
$\mu\leftarrow[\,]$\;
\For{$i\leftarrow1$ \KwTo $m$}{
  $G\leftarrow\bigcup_{j\le i}\text{goal}(S_j)$\;
  \lIf{$s$ already solved on $G$}{\textbf{continue}}
  $A\leftarrow$ single moves and verified macros that move some target of $S_i$\;
  \uIf{projection of $G$ under $A$ has $\le 10^5$ states}{
    $\nu\leftarrow$ optimal moves by greedy descent on its BFS table\;}
  \Else{
    $W\leftarrow W_0$\;
    \Repeat{$\nu$ found \textbf{or} $W > W_{\max}$}{
      $\nu\leftarrow$ \textsc{BeamSearch}$(s, A, h, W, \text{goal}=G)$\;
      $W\leftarrow4W$\;}
  }
  \lIf{no $\nu$}{\Return failure}
  $s\leftarrow s\cdot\nu$;\quad $\mu\leftarrow\mu\,\Vert\,\nu$\;
}
\Return canonicalise$(\mu)$\;
\end{algorithm}

\emph{Cost.} One evaluation is $O(K^2 d+K d^2)$ for $K=\Theta(N^2)$ cluster tokens and width $d$; a beam step costs $O(W\cdot|A|)$
evaluations. On a T4 a \Sq{10} takes about 9\,s and a \Sq{14} about 90\,s; on CPU several times longer. There are $\Theta(N^2)$
stages, so solutions have $\Theta(N^2)$ moves.

%% ---------------------------------------------------------------------
\subsection{Method 5C: net-gain scheduling of bulk commutators (new)}
\label{sec:5c}

\emph{Idea.} \Mthree applies a commutator type to a block $X\times Y$ of layer pairs only when \emph{every} cell of the
block helps the cluster it acts on. Early on many clusters need the same 3-cycle and blocks are large, but the clusters
soon diverge, the blocks shrink to single cells and batching fades: on our \Cu{20}, \MthreeSingleRoundsTwenty{} of
\Mthree's \MthreeRoundsTwenty{} center rounds fixed a single cluster. \FiveC changes only this schedule.

\emph{Steps.}
\begin{enumerate}
\item \emph{Net-gain blocks.} For a type $t$, let $g_t(i,j)$ be the change in the number of solved pieces of the cluster that
  cell $(i,j)$ acts on (it can be zero or negative), and $\lambda$ the best number of pieces per move that any
  single-cluster 3-cycle offers in the current state. \FiveC chooses
  \begin{equation}
  (X^\ast,Y^\ast)=\arg\max_{X\cap Y=\emptyset}\;\sum_{(i,j)\in X\times Y} g_t(i,j)\;-\;\lambda\,\bigl(2|X|+2|Y|+4+2|s_t|\bigr),
  \label{eq:5c}
  \end{equation}
  where $s_t$ is the type's setup and every cell of $X\times Y$ must be a verified pure 3-cycle. A block is worth applying
  exactly when its pieces per move beat the best single fix, so cells that do not help, or slightly hurt, are accepted when
  the block still pays. We maximise by alternation (rows given columns, columns given rows) from the three best rows.
\item \emph{Exact scoring.} Because several cells may act on one cluster, each candidate is \emph{applied} to the state and
  scored exactly; this also checks that nothing outside the clusters moved. Wings have one layer parameter and the same rule
  with $Y$ omitted.
\item \emph{Setup sharing.} A type whose setup equals the previous step's is charged $2|s_t|$ moves less, because the shared
  $s_t'\,s_t$ cancels when moves are canonicalised.
\item \emph{Lookahead.} With one step of lookahead (used for $N\le15$, where it is affordable), the four best candidates are
  re-scored by the best step available after them.
\item \emph{Everything else is \Mthree.} Wing parity, the \Cu{3} stage, the Square frame and edge stages, the type bank and
  the single-cluster fallback are shared, so the difference between the two methods is the schedule alone. On the Square,
  the net-gain schedule of the interior clusters (gain = drop in exact distance) and \Mthree's are both computed and the
  shorter is kept.
\end{enumerate}

\begin{algorithm}[t]
\caption{\FiveC schedule for one family of clusters (centers or wings)}
\label{alg:5c}
\small
\KwIn{state $s$; type bank (verified pure 3-cycles per type, cluster and layers); per-cluster 3-cycle libraries}
\KwOut{move sequence $\mu$ solving every cluster of the family}
$\mu\leftarrow[\,]$\;
\While{some cluster is unsolved}{
  $c_1\leftarrow$ best single-cluster 3-cycle; $\lambda\leftarrow$ pieces fixed per move of $c_1$\;
  \For{each of the 24 types with the largest total positive gain}{
    $(X,Y)\leftarrow$ alternating maximisation of Eq.~\eqref{eq:5c}\;
    apply the bulk commutator to $s$; discard it if anything outside the clusters moved\;
    record its exact gain and its moves (less shared setup)\;}
  pick the best gain per move among the blocks and $c_1$ (optionally after one step of lookahead)\;
  $s\leftarrow s\cdot\nu$;\quad $\mu\leftarrow\mu\,\Vert\,\nu$\;
}
\Return $\mu$\;
\end{algorithm}

\emph{Cost.} Each round scores 24 types, each by alternation over $O(N)$ rows and columns, plus exact application
($O(N^2)$ per candidate), so a round is $O(N^3)$ and the number of rounds is at most the number of center pieces,
$O(N^2)$; lookahead multiplies the time by about four. Measured: \FiveCTwentyMin{} minutes for a \Cu{20}. Moves: bulk
rounds fix many clusters per move, which is what moves \FiveC's ratio to $N^2/\ln N$ more slowly than the other
methods' (\S\ref{sec:scaling}); we have not proved a bound.
